% sections/methods.tex — Yöntem bölümü, taslak v1 (2026-10-02)
% Kaynaklar: CLAUDE.md (2026-10-02), 02 v2 raporu, 03 planı v2, 04_05 planı v2, 04-B planı v1,
% Tablo S-n v2, iskelet v6. Köşeli parantezli büyük harfli yerler henüz doldurulmadı.

\section{Methods}
\label{sec:methods}

Analysis identifiers in parentheses (03--07g) refer to the analysis log (table~S2) and to the code repository.

\subsection{Dataset and participants}
\label{sec:data}

We analysed the openly available dataset ``Dataset of Electroencephalograms of Juvenile Offenders'' (OpenNeuro ds006923, version 1.0.0, CC0 licence) \cite{Polo2025ds006923}. It contains resting-state EEG recordings of 140 male adolescents: 74 juvenile offenders, labelled as two subgroups (sg, $n = 49$; sg2, $n = 25$), and 66 controls without a criminal record (cg). Demographic, educational, substance-use and offence variables were taken from the dataset's participants file (table~\ref{tab:Table1}).

The dataset documentation does not define the two offender labels. A conference report from the same data-collection project describes the 74 offenders as 50 young people confined in a rehabilitation centre and 24 under non-custodial pedagogical measures \cite{RodriguezAlvarez2023}. Whether these two groups correspond to sg and sg2 cannot be verified from the published metadata. We therefore kept the dataset labels throughout.

Data collection was approved by the Ethics Committee of Universidad del Magdalena \cite{Polo2025ds006923}; for the same project, the approval date is reported as 23 May 2019, and informed consent was obtained from all participants, with the consent of parents or legal guardians for those under 18 \cite{SanchezTorres2024}. The present study is a secondary analysis of de-identified open data and required no further ethical approval.

\subsection{Recording and preprocessing by the data providers}
\label{sec:recording}

EEG was recorded with a BioSemi ActiveTwo system (128 scalp channels, 2048~Hz) \cite{Polo2025ds006923,RodriguezAlvarez2023}. The resting protocol consisted of four 60-s blocks alternating eyes closed and eyes open (closed--open--closed--open), followed by an eyes-closed block of about 8~min. The providers preprocessed the data before publication: finite impulse response band-pass filtering (1--40~Hz), average reference, artifact subspace reconstruction \cite{Mullen2015}, independent component analysis with ICLabel classification \cite{PionTonachini2019}, and downsampling to 128~Hz. Electro-oculographic and electrocardiographic channels were recorded \cite{RodriguezAlvarez2023} but are not included in the published files.

For each participant, the dataset provides the full preprocessed recording and an epochs file with four 50-s segments (two eyes closed, two eyes open). We verified that each segment is an exact excerpt of the preprocessed recording that starts 5~s after its block marker (difference 0~$\mu$V). Analyses of the four segments therefore used the epochs files, and analyses of the long eyes-closed block used the preprocessed recordings.

\subsection{Data audit and analysis samples}
\label{sec:audit}

All 140 recordings were audited before any group comparison: file integrity, duration, block markers and channel count. Five participants lacked the long eyes-closed block. Four sg participants had only about 157~s of data (one eyes-closed and one eyes-open block), and the preprocessed recording of sub-1084sg2 could not be read because its data file is not part of the published dataset. The long-block sample therefore comprised 135 participants (cg 66, sg 45, sg2 24). The alpha-reactivity analysis needs only the four 50-s segments and comprised 136 participants (cg 66, sg 45, sg2 25); the four short sg recordings entered a sensitivity analysis ($n = 140$). The sample of every analysis is given in table~\ref{tab:TableSn} and figure~\ref{fig:cohort}.

The audit also extracted the recording date and clock time of each session from the acquisition time stamps (table~\ref{tab:Table1}). The site field of the published sidecar files lists the same institution for all 140 participants, controls included; we therefore treated it as uninformative.

\subsection{Analysis log, preregistration status and use of AI tools}
\label{sec:log}

The analyses were not preregistered. Every analysis plan, decision rule and change was recorded in a time-stamped analysis log (a project file, written analysis plans and the git history). Table~S2 lists all decisions and states for each whether it was made before the relevant result was seen.

The single primary test (alpha reactivity; section~\ref{sec:primary}) was designated after the batch negative control (section~\ref{sec:negcontrol}) had been run. Before this designation, 42 unadjusted univariate comparisons of the 14 region-of-interest (ROI) features had been inspected; alpha reactivity was not among them, and these comparisons informed no feature selection (table~S2). The definitions of the last two exploratory analyses (07f, 07g) were committed to the repository in separate commits before they were run; the commit contents are provided as supplementary files. For earlier steps, the order of definition and result rests on the analysis log; for 07d and 07e, definition and results share one commit.

The analysis code was written and executed by an AI coding agent (Claude Code, Anthropic; model Claude Opus 5.5) under the authors' direction. The authors approved every analysis plan, decision rule and interpretation. Four audits of code, numbers and text were carried out in separate sessions of the same model (Claude Opus 5.5); these audits are therefore not independent. Their findings and the resulting corrections are listed in table~S2. The authors take full responsibility for the analyses and the text.

\subsection{Spectral features}
\label{sec:features}

Feature parameters were fixed before any group comparison (table~\ref{tab:TableS3}). From the long eyes-closed block, we analysed a fixed 465-s segment starting 5~s after the third eyes-closed marker, and its two halves (H1 and H2, 232.5~s each). Power spectral density was estimated with Welch's method (4-s Hann windows, 50\% overlap, 0.25-Hz resolution). Relative band power was computed for delta (1--4~Hz), theta (4--8~Hz), alpha (8--13~Hz) and beta (13--30~Hz) as a fraction of total 1--30~Hz power. The aperiodic exponent and offset were estimated with FOOOF version~1.1 \cite{Donoghue2020} in the `fixed' aperiodic mode over 3--30~Hz; a `knee' fit over 2--30~Hz was used only as a sensitivity analysis. The individual alpha frequency (IAF) was defined as the strongest FOOOF peak between 7 and 13~Hz (primary) and as the centre of gravity of the aperiodic-removed spectrum between 7 and 13~Hz (secondary).

Channels were excluded within each participant if the FOOOF fit had $R^2 < 0.9$ or if an electromyographic (EMG) index exceeded a robust $z$ score of 3 (one-sided; median $+\,3 \times 1.4826 \times$ median absolute deviation). The EMG index was the slope of log power against log frequency over 30--40~Hz; flatter slopes indicate more muscle activity. Because the providers' low-pass filter also shapes this slope, the index was used only as a relative measure. ROI features were computed from the mean spectrum of the retained channels in a posterior ROI (nine channels: [ROI CHANNEL LIST]) and in a global ROI (all retained channels).

One rule was added after inspection of quality-control data but before any group comparison: participants with more than 25\% of channels excluded were flagged as having an atypical spectrum, and only the EMG rule was applied to them. Two sg2 participants were flagged (88 and 57 channels excluded by the $R^2$ rule). They were kept in the primary analyses and excluded in sensitivity analyses. When the rule was adopted, it was known that both belonged to sg2 and used cannabis and cocaine (table~S2).

Vigilance was indexed by the slope of the log$_{10}$ theta/alpha ratio across consecutive 30-s windows of the long block; a positive slope indicates increasing drowsiness \cite{Strijkstra2003}. The 14 ROI features used in the ROI-level analyses were relative delta, theta, alpha and beta power, the aperiodic exponent and offset, and the primary IAF, each in the posterior and global ROIs.

\subsection{Batch negative control (03)}
\label{sec:negcontrol}

Before any offender--control comparison, we asked (a) whether the 14 ROI features separate the two offender subgroups (sg vs sg2) and (b) whether a model trained on sg vs cg also separates sg2 from cg. A preliminary comparison showed that sg and sg2 differ not only in recording period but also in offence profile, alcohol use and age (table~\ref{tab:Table1}). A positive result in (a) could therefore not be read as a pure batch effect.

The pipeline was median imputation, standardisation and L2-regularised logistic regression with balanced class weights, all fitted within training folds. Each participant contributed one row. The outer loop was stratified 5-fold cross-validation repeated 20 times; an inner stratified 5-fold loop selected the regularisation strength (13 values from $10^{-3}$ to $10^{3}$) by area under the ROC curve (AUC). In (b), the model of each outer fold scored that fold's held-out controls and all 24 sg2 participants, and the AUC of sg2 versus the held-out controls was computed within the fold and averaged over the 100 folds. For both tests, labels were permuted 1000 times and the full nested procedure was rerun with 5 repeats per permutation; $p = (k+1)/(N+1)$ and $\alpha = 0.025$ per test (Bonferroni). Six prespecified sensitivity analyses repeated the pipeline without permutation: excluding the two atypical participants; regressing the EMG index and the vigilance slope, or age, out of each feature within training folds; using H1 instead of the full segment; omitting the aperiodic offset; and using the knee-model features.

\subsection{Leakage demonstration on the published features (04-A, 04-B)}
\label{sec:leakage}

The dataset also contains precomputed features: 2048 spreadsheet files (four bands $\times$ four segments [C1, O1, C2, O2] $\times$ 128 channels), each with seven statistics per participant, and an equivalent MATLAB file. We first audited these files (04-A): consistency across files, column quality, internal consistency of the statistics, and the identity of the 112 participants they contain (cg 66, sg 46; no sg2). Identity was checked by correlating each participant's standardised log-power profile in the spreadsheets with the same profile recomputed from the published epochs files (table~\ref{tab:TableS5}).

We then quantified leakage with these features as published (04-B): 448 rows (112 participants $\times$ 4 segments) with 3584 features each. A $2 \times 2$ design crossed the cross-validation unit (row-wise stratified 10-fold vs subject-wise 5-fold) with the place of feature selection (ANOVA $F$ test, top 100 features, on all data vs within training folds). This gives four non-nested pipelines, N1--N4, each repeated with 10 random seeds; N1, with row-wise splits and selection on all data, represents a common naive approach. We built these naive pipelines to illustrate the error; they do not reproduce any published analysis of these data. The correct pipeline (P) was nested and subject-wise: an outer stratified group 5-fold loop repeated 20 times and an inner stratified group 5-fold loop that selected the number of features ($k \in \{10, 50, 100, 500\}$) and, for the support vector machine, $C \in \{0.1, 1, 10\}$. Classifiers were a support vector machine with a radial basis function kernel and a random forest with 500 trees; standardisation was always fitted within training folds. A participant's score was the mean of its row scores. Subject-level AUC was the primary metric; row-level AUC, balanced accuracy and Golden Distance were secondary.

To show how much performance arises without signal, labels were permuted at the subject level 200 times for N1 (both classifiers) and for P (support vector machine, 5 repeats per permutation). Sensitivity analyses used only the 100 identity-verified participants, removed the kurtosis features, or used only the eyes-closed segments.

\subsection{Primary hypothesis test: alpha reactivity (05)}
\label{sec:primary}

The single primary test compared the posterior alpha reactivity index (ARI) between offenders (sg and sg2 combined) and controls. Absolute alpha power (8--13~Hz; Welch's method as above) was computed in each of the four 50-s segments of the first 4~min and averaged over the nine posterior channels, omitting each participant's excluded channels (section~\ref{sec:features}; all nine channels for sub-1084sg2). With $C$ and $O$ the mean alpha power of the two eyes-closed and the two eyes-open segments, $\mathrm{ARI} = (C - O)/(C + O)$. Because the index is a ratio, multiplicative gain differences cancel.

The test was a two-sided Mann--Whitney $U$ test ($\alpha = 0.05$; 70 offenders, 66 controls). Effect sizes were the rank-biserial correlation and the Hodges--Lehmann shift \cite{HodgesLehmann1963}, each with a bootstrap 95\% confidence interval (CI; 2000 resamples). In a covariate sensitivity analysis, ordinary least-squares models regressed the ARI, and separately its ranks, on group, age and the posterior EMG index from the same segments, with HC3 standard errors \cite{MacKinnonWhite1985}. A second sensitivity analysis added the four short sg recordings, using their single eyes-closed and eyes-open segments ($n = 140$). The interpretation was fixed in advance: a positive result would be reported as a group difference not explained by gain but not separable from substance use, recording period or site; a negative result would be quantified by its CIs.

\subsection{Quantifying the ROI-level null result (06)}
\label{sec:null}

Offenders (sg and sg2 combined, $n = 69$) and controls ($n = 66$) were classified with the 14 ROI features and the pipeline of section~\ref{sec:negcontrol}. Each participant's out-of-fold probability was averaged over the 20 repeats. The AUC of these averaged scores received a 95\% CI from a class-stratified participant bootstrap (2000 resamples, percentile) and, as a cross-check, from DeLong's method \cite{DeLong1988}. This CI does not include variability from model training; the 2.5th--97.5th percentile range of the per-repeat AUCs is reported separately. The upper CI limit was taken as the smallest effect that is approximately excluded and was converted to Cohen's $d = \sqrt{2}\,\Phi^{-1}(\mathrm{AUC})$.

\subsection{Exploratory channel-level analyses (07--07g)}
\label{sec:exploratory}

After the ROI-level analyses, a series of exploratory analyses asked whether channel-level spectral features separate sg from controls and whether such a separation transfers to sg2. For comparability with the leakage analysis, these analyses used the 100 identity-verified participants of the spreadsheet cohort (sg 45, cg 55), with features computed from the published recordings. Two feature sets were computed per segment and channel (Welch's method as above; 128 channels $\times$ 4 bands = 512 features): (A) log$_{10}$ absolute band power and (B) relative band power. All analyses used pipeline P with the support vector machine (section~\ref{sec:leakage}) and 200 subject-level label permutations with 5 repeats each. Their $p$ values are not adjusted for multiple comparisons.

\begin{itemize}
  \item 07: sets A and B on all four segments. Separation in A but not in B was to be attributed to absolute amplitude.
  \item 07a: set B on the eyes-closed and the eyes-open segments separately. 07b: set B on the eyes-closed segments without the 32 channels of the BioSemi C block, which covers frontopolar, anterior frontal and right frontal sites (equivalents include Fp1, Fpz, Fp2, AF7, AF3, AFz, AF4, AF8, F1, Fz, F2, F4, F6, F8, FC1, FCz and FC2).
  \item 07c and 07d: transfer to sg2 ($n = 24$) of the eyes-closed model (07c) and of the three strongest segment models (set A on all segments, set B on all segments and set B on the eyes-open segments; 07d). Transfer was measured with the within-fold AUC of section~\ref{sec:negcontrol}; the permutation null shuffled the training labels.
  \item 07e: the long eyes-closed block (the 465-s segment of section~\ref{sec:features}) divided into four consecutive, non-overlapping 116.25-s parts, set B, without channel exclusion. 07f: transfer of the 07e model to the long blocks of sg2 ($n = 24$).
\end{itemize}

The decision rule of 07a--07c combined an AUC threshold with $p < 0.05$ and left the interval 0.60--0.65 undefined; 07d--07f used a single criterion, permutation $p < 0.05$. For the transfer tests, $p < 0.05$ in any model would have narrowed the claim of this paper; without multiplicity correction, this rule is conservative for the claim that there is no evidence of transfer.

Analysis 07g quantified uncertainty without a decision rule. For the six source models (07 A, 07 B, 07a eyes closed, 07a eyes open, 07b and 07e), each participant's decision value was averaged over the 20 repeats, and the AUC of these averaged scores received a 95\% CI from a class-stratified participant bootstrap (2000 resamples). For the five transfer models (07c, three 07d models and 07f), a weighted participant bootstrap drew multinomial weights for the controls and for sg2, shared them across all folds of a resample, recomputed each fold's AUC as a weighted Mann--Whitney statistic and averaged over the 100 folds (2000 resamples, percentile). Before the CIs were computed, all 1120 reproduced AUCs were checked against the saved values (maximum difference $6 \times 10^{-17}$). These CIs reflect participant sampling only. A descriptive map of the 07e contrast (Hedges' $g$ for each channel and band; no test) is shown in figure~\ref{fig:S07e}.

\subsection{Statistics, metrics and software}
\label{sec:stats}

The AUC was the primary classification metric; balanced accuracy and the Golden Distance (GD) were secondary. GD summarises three or more performance metrics in a single score \cite{Melek2025GD}. We used accuracy, sensitivity, specificity and F1 with the offender class as positive, for which the GD polygon is a square; lower values indicate performance closer to the ideal. Permutation $p$ values were computed as $(k+1)/(N+1)$ \cite{Ojala2010}; with 200 permutations the smallest attainable value is $1/201$, reported as $p \le 0.005$. Ranges across cross-validation repeats (2.5th--97.5th percentile) describe training variability and are not confidence intervals. Exploratory $p$ values are not adjusted for multiple comparisons. AUCs and CIs are given with two decimals in the text, because the bootstrap limits vary by about $\pm 0.01$ with the random seed, and $p$ values with three decimals. Random seeds were fixed (20261001).

Analyses were run in Python with scikit-learn \cite{Pedregosa2011}, FOOOF \cite{Donoghue2020} and [OTHER PACKAGES AND VERSIONS]. The code, the analysis log and all result files are available at [REPOSITORY / ZENODO DOI].
